\clearpage
\appendix
\section*{Appendices}
\section{Open-term typing (proposal)}

% (Section title supplied by \section above.)

Python uses a form of \emph{late binding}, where most variables only need to be defined at execution time. For example, 
the following program is valid: 

\begin{python}
def foo():
    return x

x = 10
y = foo()
\end{python}

While \pyinline{x} is not defined before its usage inside the definition of \pyinline{foo}, it is defined by the
time the function is called. Therefore, Python uses a limited form of \emph{dynamic scoping}. In PurePy, we would
like to restrict this form of scoping to something that is easier to statically analyse, and less error-prone 
(see Section~\ref{sec:captured-vars}). However, some forms of late binding must be allowed, as we could otherwise not
define (mutually) recursive functions:

\begin{python}
def even(n):
    if n == 0: return True
    else: odd (n-1)

def odd(n):
    if n == 0: return False
    else: return even (n-1)
\end{python}

The function \pyinline{even} refers to \pyinline{odd} before that function is defined. We restrict the
allowed forms of late binding by the introduction of a type system, which uses \emph{open term types},
i.e., a term is allowed to be dependent on yet undefined variables.

Figure~\ref{fig:open-term-types} introduces several auxiliary symbols we need for the typing rules. Each
open term (i.e, statement or expression) has an \emph{open term type} $\tau$, which denotes the set of
variable names that are unbound in the respective term.

As with other programming languages, scoping in Python and PurePy has nesting, i.e., each \pyinline{def}
and \pyinline{lambda} opens a new scope. Therefore, the typing rules are expressed in terms of 
\emph{frames} $\Psi$, which correspond to the different nesting levels. A particular frame $\Psi$ consists
of the set of \emph{closure variables} $\Sigma$, i.e., the set of variables defined in the current frame
that are used within a closure (see Section~\ref{sec:captured-vars}), and an \emph{environment} $\Gamma$,
which maps variable names to open term types. Finally, a \emph{frame stack} $\Delta$ holds the nested
scoping frames.

Figures~\ref{fig:open-stmts-type-rules} and \ref{fig:open-expr-type-rules} illustrate the typing relation for statements
and expressions respectively. Both have the same shape, where a statement or an expression evaluates under a
stack $\Delta$ to an open type $\tau$ and an updated stack $\Delta'$. Statements can change the stack by introducing
new bindings in $\Gamma$ through assignments at the top-most frame, while expressions can change the set of closure 
variables $\Sigma$ at any nesting depth (e.g., a free variable might be closed by a binding which is defined at any 
enclosing scope).

We will start by illustrating the typing rules for expressions. A constant has an empty type, and does not change
the stack. For variables, there are three different cases. In $\ruleName{var1}$, the variable is
already bound in the current scope $\Gamma$. Therefore, its type can be read from the environment, and the
variable itself is added to the set of closure variables at the current scope level (to circumvent the issue explained in
Section~\ref{sec:captured-vars}). If the variable is not bound in the current environment, it is typed at the
surrounding scope as shown in $\ruleName{var2}$. The base-case is illustrated by $\ruleName{var3}$, where a variable
is typed at the global scope (i.e. $\carrow{0}$). If the variable is not defined at the global scope either, it receives
itself as an open term type.

Lambdas introduce a new frame onto the stack, where the set of closure variables is empty, and the environment only
maps the bound parameters of the lambda to empty open types. The body of the lambda term is then typed at a nested
scope (i.e., $\carrow{n+1}$).

For function application, two different rules apply. If the function application happens at global scope (i.e.,
$\carrow{0}$), then all terms of the application must have empty open types (i.e., they cannot depend on variables
that are not yet defined). If the application is inside a nested scope (e.g., under a \pyinline{def} or \pyinline{lambda})
then each individual term is typed in sequence, where the returned type is then the union of the open term types, and
the returned stack the one of the last typing judgement.

For statements, the $\ruleName{pass}$ is trivial. For an assignment, the assigned variable cannot be part of the set of
closure variables at the current scoping level. The returned type is the type of the right-hand side expression without
the name currently being defined. The function $\mathrm{close}$ removes the second argument from all bindings in the
updated environment (i.e., $\mathrm{close}(\Gamma, x)(y) = \Gamma(y) \setminus \{x\}$).

The rules for conditional evaluation ($\ruleName{if1}$ and $\ruleName{if2}$) follow a similar structure to the ones for
function evaluation. If the conditional statement is on the global scope (i.e., $\carrow{0}$), then all sub-expressions
and sub-statements must be closed. If it is inside a nested scope, each sub-term is typed in sequence, and the
returned type is again the union of the open term types of the sub-terms.  

The rule for function definition $\ruleName{def}$ works similarly to the one for lambda expressions, with the added step of
closing the environment by the function name $f$ in the returned environment stack.

The rules for return and empty return are trivial again, they have a restriction on the nesting depth, so that they can only
appear inside function definitions and not on a global scope.

The rule for statement sequence is also trivial.

\begin{figure}
\begin{subfigure}{\textwidth}
\small
\begin{center}
\begin{tabular}{rcll}
$\tau$ & $\subseteq$ & $\mathbb{V}$ & (open term type) \\
$\Psi$ & $::=$ & $\Sigma \times \Gamma$ & (frame) \\
$\Sigma$ & $::=$ & $\{ x \}$ & (closure variables) \\
$\Gamma$ & $::=$ & $\mathbb{V} \rightarrow \mathcal{P}(\mathbb{V})$ & (environment) \\
$\Delta$ & $::=$ & $\Delta : \Psi$ & (frame stack) \\
 & $\mid$ & $\varnothing$ & (empty stack)
\end{tabular}
\end{center}
\end{subfigure}
\caption{Open-term types}
\label{fig:open-term-types}

\begin{subfigure}{\textwidth}
\flushleft \shadebox{$\Delta \vdash s \carrow{n} \tau, \Delta'$}
\begin{mathpar}
\inferrule*[lab={\ruleName{pass}}]
{
    \strut
}
{
    \Delta \vdash \exPass \carrow{n} \Delta
}
\and%
\inferrule*[lab={\ruleName{assign}}]
{
    x \notin \Sigma \qquad \Delta : (\Sigma, \Gamma) \vdash e \carrow{n} \tau, \Delta' : (\Sigma', \Gamma')  
}
{
    \Delta : (\Sigma, \Gamma) \vdash \exAssign{x}{e} \carrow{n} \gamma \setminus \{ x \},  
    \Delta' : (\Sigma', \mathrm{close}(\Gamma'[x \mapsto \gamma], x)) 
}
\and%
\inferrule*[lab={\ruleName{def}}]
{
    \Delta : (\varnothing, [f \mapsto \varnothing, \overline{x} \mapsto \varnothing]) \vdash s \carrow{n+1} \tau, 
    \Delta' : (\Sigma', \Gamma') : \underline{\;\;}
}
{
    \Delta \vdash \exDef{f}{\bar{x}}{s} \carrow{n} \tau, \Delta' : (\Sigma', \mathrm{close}(\Gamma'[f \mapsto \tau], f))
}
\and%
\inferrule*[lab={\ruleName{return}}]
{
    n > 0 \qquad \Delta \vdash e \carrow{n} \tau, \Delta'
}
{
    \Delta \vdash \exReturn{e} \carrow{n} \tau, \Delta'
}
\and%
\inferrule*[lab={\ruleName{returnNone}}]
{
    n > 0
}
{
    \Delta \vdash \exReturnNone \carrow{n} \varnothing, \Delta
}
\and%
\inferrule*[lab={\ruleName{seq}}]
{
    \Delta \vdash s_1 \carrow{n} \tau_1, \Delta' \qquad \Delta' \vdash s_2 \carrow{n} \tau_2, \Delta''
}
{
    \Delta \vdash \exSeq{s_1}{s_2} \carrow{n} \tau_1 \cup \tau_2, \Delta''
}
\end{mathpar}
\end{subfigure}
\caption{Open statements typing rules}
\label{fig:open-stmts-type-rules}

\begin{subfigure}{\textwidth}
\flushleft \shadebox{$\Delta \vdash e \carrow{n} \tau, \Delta'$}
\begin{mathpar}
\small
\inferrule*[lab={\ruleName{const}}]
{
    \strut
}
{
    \Delta \vdash n \carrow{n} \varnothing, \Delta
}
\and%
\inferrule*[lab={\ruleName{var1}}]
{
    x \in \mathrm{dom}(\Gamma)
}
{
    \Delta : (\Sigma, \Gamma) \vdash x \carrow{n} \Gamma (x), \Delta : (\Sigma \cup \{ x \}, \Gamma)
}
\and%
\inferrule*[lab={\ruleName{var2}}]
{
    x \notin \mathrm{dom}(\Gamma) \qquad \Delta \vdash x \carrow{n-1} \tau, \Delta'
}
{
    \Delta : \Psi \vdash x \carrow{n} \tau, \Delta' : \Psi
}
\and%
\inferrule*[lab={\ruleName{var3}}]
{
    x \notin \mathrm{dom}(\Gamma)
}
{
    \varnothing : (\Sigma, \Gamma) \vdash x \carrow{0} \{x\}, \varnothing : (\Sigma, \Gamma) 
}
\and%
\inferrule*[lab={\ruleName{lambda}}]
{
    \Delta : (\varnothing, [\overline{x} \mapsto \varnothing]) \vdash e \carrow{n+1} \tau, \Delta'
}
{
    \Delta \vdash \exLambda{\overline{x}}{e} \carrow{n} \tau, \Delta'
}
\and%
\inferrule*[lab={\ruleName{funcApp1}}]
{
    \forall i: \Delta \vdash e_i \carrow{0} \varnothing, \Delta
}
{
    \Delta \vdash e_1 (e_2, \dots, c_n) \carrow{0} \varnothing, \Delta
}
\and%
\inferrule*[lab={\ruleName{funcApp2}}]
{
    n > 0 \qquad \Delta \vdash e_1 \carrow{n} \gamma_1, \Delta^{[1]} \qquad \Delta^{[1]} \vdash e_2 \carrow{n} \gamma_2, 
    \Delta^{[2]} \qquad \cdots
}
{
    \Delta \vdash e_1 (e_2, \dots, e_n) \carrow{n}  \bigcup \gamma_i, \Delta^{[n]}
}
\end{mathpar}
\end{subfigure}
\caption{Open expressions typing rules}
\label{fig:open-expr-type-rules}
\end{figure}

